\section{How Does My Paper Relate to the ACC/IDS Problem?}

\begin{contextbox}{Background: The Chain of Papers}
	The existing research establishes:
	\begin{enumerate}[label=\textbf{\arabic*.}]
		\item \textbf{Adaptive Cruise Control (ACC)} regulates vehicle speed and
		      following distance using a PID controller. The PID computes a
		      control signal from the speed error:
		      \[
			      u(t) = K_p\, e(t) + K_i \int_0^t e(\tau)\,d\tau + K_d \frac{de(t)}{dt}
		      \]
		      where $e(t) = S_t(t) - S_c(t)$ is the difference between target speed
		      and current (sensor-reported) speed.
		      
		\item \textbf{The Intrusion Detection System (IDS)} is a passive
		      extension to the ACC. It monitors the \textbf{Control Area Network(CAN)}
		      bus for spoofed speed messages using machine learning. When it
		      detects an attack, it raises a flag.
		      
		\item \textbf{Paper \cite{jedh2023improvement}} investigates the impact of \textbf{cold braking}
		      as a mitigation strategy.
	\end{enumerate}
\end{contextbox}

\begin{questionbox}{Question: What does my paper do differently?}
	My paper investigates an alternative mitigation approach.
	Instead of cold braking when the IDS detects an attack, I propose:
	
	\begin{enumerate}[label=(\alph*)]
		\item Using the \textbf{tracking control equation}\textsubscript{\cite{tymoshchuk2024tracking}} to continuously track the error between the expected speed (what the vehicle should be doing based on the reference trajectory) and the reported speed (what the potentially-spoofed sensor claims).
		      
		\item \textbf{Correcting for the error} using the designed control function $u(k)$\textsubscript{\cite{tymoshchuk2024tracking}},
		      which computes the input needed to drive the system state back toward the reference trajectory.
		      
		\item Performing a \textbf{mode switch} if the tracking error exceeds a threshold,
		      i.e. if the discrepancy between expected and reported speed grows too large,
		      the system can escalate its response (e.g., increase correction aggressiveness
		      or fall back to emergency braking).
	\end{enumerate}
	
	Additionally, I \textbf{simulate the attack scenarios} to demonstrate that the tracking controller successfully corrects for injected speed errors, showing convergence of the tracking error to zero (or near-zero) even under adversarial perturbation.
\end{questionbox}

\begin{contextbox}{Summary of Positioning}
	\begin{center}
		\renewcommand{\arraystretch}{1.4}
		\begin{tabular}{|l|l|l|}
			\hline
			\textbf{Component} & \textbf{Paper \cite{jedh2023improvement} Approach} & \textbf{My Approach}               \\
			\hline
			Detection          & ML-based IDS on CAN bus                            & Same (IDS flags the attack)        \\
			\hline
			Mitigation         & Cold brake (binary)                                & Tracking control correction        \\
			\hline
			Error model        & PID speed error $e(t)$                             & Discrete-time $e(k) = r(k) - x(k)$ \\
			\hline
			Correction         & None (just stop)                                   & $u(k)$ drives $e(k) \to 0$         \\
			\hline
			Escalation         & N/A                                                & Mode switch if $\|e(k)\| > \theta$ \\
			\hline
		\end{tabular}
	\end{center}
\end{contextbox}

\bibliographystyle{ieeetr}

\bibliography{references}
